% Part: turing-machines
% Chapter: undecidability
% Section: unsolvability-decision-problem

\documentclass[../../../include/open-logic-section]{subfiles}

\begin{document}

\olfileid{tur}{und}{dec}
\olsection{فیصلے دا مسئلہ}

اسی پہلے درجے دی منطق نوں \emph{قابلِ فیصلہ} اُدوں تے
صرف اُدوں کہندے آں جدوں کوئی مؤثر طریقہ ایہ طے کردا
ہووے کہ دتا ہویا !!{sentence} درست اے یا نہیں۔ پر ایہو
جیہا طریقہ نہیں: پہلے درجے دے جملیاں دی درستی دا فیصلہ کرن
والا مسئلہ حل نہیں ہو سکدا۔

ایہ اہم منفی نتیجہ قائم کرن لئی اسی ثابت کراں گے کہ
ٹیورنگ مشین فیصلے دا مسئلہ حل نہیں کر سکدی۔ یعنی
ایہو جیہی کوئی ٹیورنگ مشین نہیں جیڑی، جدوں وی اوہ
ایس پٹی اُتے شروع ہووے جیہدے اُتے پہلے درجے دا !!{sentence}
لکھیا ہووے، آخرکار رُکے تے ایہ ویکھ کے کہ اوہ !!{sentence}
درست اے یا نہیں، $1$ یا~$0$ آؤٹ پٹ دے۔ چرچ--ٹیورنگ دے
دعوے مطابق ہر حساب کرن جوگ فنکشن ٹیورنگ مشین نال
وی حساب ہو سکدا اے۔ ایس لئی جے ایہ “درستی دا
فنکشن” کسے مؤثر طریقے نال حساب ہو سکدا ہووے،
تاں ٹیورنگ مشین نال وی حساب ہو سکدا اے۔ جے ایہ
ٹیورنگ مشین نال حساب نہیں ہو سکدا، تاں کوئی مؤثر
طریقہ وی ایہنوں حساب نہیں کر سکدا۔

فیصلے دے مسئلے دا ناقابلِ حل ہونا ثابت کرن لئی
اسی رُکّن دے مسئلے نوں فیصلے دے مسئلے تک گھٹاواں گے۔
ایس دا مطلب ایہ اے: اسی ثابت کر چکّے آں کہ فنکشن~$h(e,w)$،
جیہدی قدر $1$ اے جے~$e$ دے بیان والی ٹیورنگ مشین
انپٹ~$w$ اُتے رُکے تے دوجی صورت وچ~$0$ اے،
ٹیورنگ مشین نال حساب نہیں ہو سکدا۔ اسی دکھاواں گے کہ
جے پہلے درجے دے جملیاں دی درستی دا فیصلہ کرن والی ٹیورنگ مشین ہوندی،
تاں~$h$ دا حساب کرن والی ٹیورنگ مشین وی ہوندی۔
چونکہ $h$ ٹیورنگ مشین نال حساب نہیں ہو سکدا،
ایس لئی درستی دا فیصلہ کرن والی ٹیورنگ مشین
وی نہیں ہو سکدی۔

ایس حکمتِ عملی دا پہلا قدم ایہ دکھانا اے کہ ہر
انپٹ~$w$ تے ٹیورنگ مشین~$M$ لئی اسی مؤثر
طریقے نال اک !!a{sentence} $!T(M, w)$ بیان
کر سکدے آں، جیڑا~$M$ دیاں ہدایات تے
انپٹ~$w$ نوں ظاہر کرے، تے اک !!a{sentence}~$!E(M, w)$
وی بیان کر سکدے آں، جیڑا ایہ آکھے کہ “$M$
آخرکار رُکدی اے”، ایس طرح کہ:
\begin{quote}
  $\Entails !T(M, w) \lif !E(M,w)$ اُدوں تے صرف اُدوں جدوں $M$
  انپٹ~$w$ اُتے رُکدی اے۔
\end{quote}
ساڈے ثبوت دا وڈا حصہ ایہ جملے $!T(M, w)$
تے~$!E(M, w)$ بیان کرن تے ایہ جانچن وچ
لگے گا کہ $!T(M, w) \lif !E(M, w)$ اُدوں تے صرف
اُدوں درست اے جدوں $M$~انپٹ~$w$ اُتے رُکدی اے۔

\end{document}
